\documentclass{article}
\usepackage[preprint]{neurips_2025}
\usepackage[utf8]{inputenc}
\usepackage[T1]{fontenc}
\usepackage{hyperref}
\usepackage{url}
\usepackage{booktabs}
\usepackage{amsfonts}
\usepackage{nicefrac}
\usepackage{microtype}
\usepackage{graphicx}
\usepackage{natbib}

\title{Estrogen Receptor-Negative, Progesterone Receptor-Positive, HER2-Negative Breast Cancer: A Systematic Review of Adjuvant Endocrine Therapy}

\author{AutoResearchClaw}

\begin{document}
\maketitle


\begin{abstract}


Estrogen receptor-negative/progesterone receptor-positive/HER2-negative (ER−/PR+/HER2−) breast cancer is a rare and biologically controversial subtype accounting for approximately 1–4% of all invasive breast cancers. The presence of progesterone receptor expression in the absence of estrogen receptor activity challenges the canonical understanding of hormone receptor biology, wherein PR transcription is classically driven by ER-mediated signaling. Whether this phenotype represents a genuine biological entity or a technical artifact of immunohistochemical (IHC) assessment remains actively debated. Critically, the role of adjuvant endocrine therapy in this subtype is uncertain, as landmark trials and meta-analyses—including the Early Breast Cancer Trialists' Collaborative Group (EBCTCG) overview—have consistently identified ER status, not PR status, as the principal predictor of endocrine therapy benefit. This review synthesizes the available evidence on the biology, clinical behavior, diagnostic challenges, and therapeutic implications of ER−/PR+/HER2− breast cancer, with particular focus on the role of adjuvant endocrine therapy. We examine data from large database analyses, the EBCTCG meta-analysis, molecular profiling studies, and current guideline recommendations to provide a comprehensive assessment of this clinically challenging subtype and to propose a framework for clinical decision-making.

\end{abstract}

\section{Introduction}

Breast cancer is a molecularly heterogeneous disease, and the classification of tumors according to hormone receptor (HR) and human epidermal growth factor receptor 2 (HER2) status has been fundamental to treatment selection for decades \citep{perou2000}. Among the four canonical immunohistochemical subtypes (HR+/HER2−, HR+/HER2+, HR−/HER2+, and triple-negative), the estrogen receptor-negative/progesterone receptor-positive (ER−/PR+) phenotype occupies a peculiar and contested position.

The progesterone receptor gene (PGR) is a well-characterized estrogen-responsive gene; its transcription is largely regulated by activated estrogen receptor-alpha (ERα) \citep{horwitz1978}. Consequently, the expression of PR in the absence of ER is biologically counterintuitive, and has led many investigators to question whether ER−/PR+ tumors represent a true biological entity or are primarily the product of pre-analytical, analytical, or interpretive errors in immunohistochemical testing \citep{kunc2018, hefti2014}.

The clinical significance of this question extends directly to treatment: adjuvant endocrine therapy is a cornerstone of breast cancer management for HR-positive disease, reducing recurrence and mortality through estrogen deprivation (aromatase inhibitors) or receptor antagonism (tamoxifen) \citep{ebctcg2011}. However, if ER−/PR+ tumors lack functional ER signaling, the biological rationale for endocrine therapy is uncertain. If, conversely, these tumors harbor cryptic ER activity or alternative mechanisms of hormone sensitivity, then withholding endocrine therapy may deprive patients of potential benefit.

This review aims to: (1) examine the biological basis and diagnostic challenges of the ER−/PR+ phenotype; (2) evaluate the clinical behavior and prognosis of ER−/PR+/HER2− breast cancer; (3) critically assess the evidence for and against adjuvant endocrine therapy in this subtype; and (4) propose a framework for clinical decision-making based on current evidence and guideline recommendations.

\section{Biology of the ER−/PR+ Phenotype}

\subsection{Classical Hormone Receptor Signaling}

In normal breast epithelium and in most hormone receptor-positive breast cancers, estrogen binds to ERα, which then dimerizes and translocates to the nucleus, where it acts as a transcription factor for estrogen-responsive genes—including the progesterone receptor gene (PGR) \citep{horwitz1978}. This canonical pathway predicts that PR expression should be contingent upon functional ER activity, making the ER−/PR+ phenotype paradoxical.

\subsection{ER-Independent Mechanisms of PR Expression}

Several mechanisms have been proposed to explain PR expression in the absence of detectable ER:

First, alternative transcription factor activation may drive PGR expression independently of ER. Growth factor signaling through the PI3K/AKT/mTOR pathway, MAPK pathway, or other kinase cascades can activate PGR transcription through ER-independent mechanisms \citep{cui2005}. The RANKL/RANK/NF-κB axis has also been shown to upregulate PR expression in certain contexts \citep{gonzalez-suarez2010}.

Second, epigenetic mechanisms including promoter methylation or histone modification of the ESR1 gene may silence ER expression while leaving PGR expression intact, particularly if PGR regulation has become uncoupled from ER control through chromatin remodeling \citep{lapidus1996}.

Third, ER splice variants or mutations may render ER undetectable by standard IHC antibodies while retaining partial transcriptional activity sufficient to drive PR expression. Conformational changes in the ER protein could prevent antibody binding without abolishing receptor function \citep{kunc2018}.

\subsection{The Artifact Hypothesis}

A substantial body of evidence suggests that many, perhaps most, ER−/PR+ results arise from technical limitations of immunohistochemistry:

False-negative ER results may occur due to prolonged cold ischemia time, suboptimal fixation, antigen retrieval failure, or antibody-specific technical issues \citep{allison2020}. Tumor heterogeneity—particularly in mixed tumors with both ER-positive and ER-negative components—may lead to discordant results depending on which area is sampled or predominates in the tissue section.

False-positive PR results may arise from antibody cross-reactivity, nonspecific staining, or interpretive error, particularly at low expression levels \citep{hefti2014}.

Hefti et al. \citep{hefti2014} analyzed gene expression data and found that PGR expression shows minimal variation across ER-negative breast cancers—less variable than 99% of genes in the genome in the ER-negative context. This finding strongly suggests that apparent PR positivity in ER-negative tumors is unlikely to carry biological or clinical significance, and that the ER−/PR+ phenotype as detected by IHC is not a reproducible molecular subtype.

Supporting this, as immunohistochemical methods have improved over the past three decades, the reported frequency of the ER−/PR+ phenotype has decreased from approximately 4% in the early 1990s to approximately 1% in recent large database studies \citep{kunc2018}, consistent with a declining rate of technical artifact.

\section{Epidemiology and Clinical Characteristics}

\subsection{Prevalence}

The ER−/PR+ phenotype accounts for approximately 1–4% of all invasive breast cancers across large series, making it the rarest of the four HR/HER2 subtype combinations \citep{kunc2018, deyarmin2015}. In a large analysis from the National Cancer Database (NCDB), approximately 2% of breast cancers were classified as ER−/PR+ \citep{asco2021ncdb}.

\subsection{Clinicopathological Features}

When compared with ER+/PR+ tumors, ER−/PR+ breast cancers consistently demonstrate more aggressive features:

Regarding age at diagnosis, ER−/PR+ tumors disproportionately affect younger women, with a median age of approximately 50 years compared to 59 years for ER+ tumors \citep{deyarmin2015}.

Regarding tumor grade, the majority of ER−/PR+ tumors are high-grade (grade 3), with one NCDB analysis reporting 80% grade 3, 18% grade 2, and only 2% grade 1 \citep{rala2023}.

Regarding proliferation, Ki-67 indices tend to be higher in ER−/PR+ tumors than in ER+/PR+ tumors, and are more comparable to triple-negative breast cancer \citep{kunc2018}.

Regarding molecular subtype distribution, when profiled by gene expression assays such as PAM50, ER−/PR+ tumors are heterogeneous. While many classify as basal-like or HER2-enriched (similar to triple-negative disease), a notable proportion—up to 25% of clinically ER-negative tumors in one large PAM50 study—classify as luminal \citep{prat2012}. This heterogeneity underscores the likely mixed etiology of the ER−/PR+ designation.

\subsection{Prognosis}

The prognosis of ER−/PR+/HER2− breast cancer is intermediate, falling between ER+/PR+ (most favorable) and triple-negative (least favorable) subtypes. Deyarmin et al. \citep{deyarmin2015} reported that ER−/PR+ tumors demonstrated worse overall survival than ER+ tumors but better outcomes than ER−/PR− tumors, suggesting some degree of residual hormone sensitivity or a biologically less aggressive phenotype relative to triple-negative disease.

However, the heterogeneity within this subtype—particularly when stratified by HER2 status—must be emphasized. The NCDB analysis presented at ASCO 2021 \citep{asco2021ncdb} found that ER−/PR+/HER2+ tumors may behave differently from ER−/PR+/HER2− tumors, with the former potentially deriving benefit from endocrine therapy while the latter did not. This review focuses specifically on the HER2-negative subset.

\section{Evidence for Adjuvant Endocrine Therapy}

\subsection{The EBCTCG Meta-Analysis}

The most authoritative evidence on adjuvant endocrine therapy in breast cancer comes from the Early Breast Cancer Trialists' Collaborative Group (EBCTCG) patient-level meta-analyses. The 2011 EBCTCG overview \citep{ebctcg2011} analyzed data from 20 randomized trials encompassing 21,457 women treated with approximately 5 years of adjuvant tamoxifen.

The central finding was unequivocal: in ER-positive disease, tamoxifen substantially reduced recurrence (rate ratio 0.53 during years 0–4; 0.68 during years 5–9) and breast cancer mortality (rate ratio 0.71 during years 0–4). In ER-negative disease, however, tamoxifen had little or no effect on recurrence or mortality.

Critically for the ER−/PR+ question, the EBCTCG analysis found that PR status was not independently predictive of tamoxifen benefit. In ER-positive disease, the proportional reduction in recurrence was approximately the same regardless of PR status. Among ER-negative tumors, including those that were PR-positive, the overall effects of tamoxifen were small and not statistically significant.

The meta-analysis also identified a relatively sharp threshold in tamoxifen efficacy with respect to quantitative ER measurement: minimal benefit was observed at very low ER levels (4–9 fmol/mg by ligand-binding assay), with substantial benefit emerging at higher levels (≥10 fmol/mg). However, the confidence intervals at very low ER levels were wide, precluding definitive conclusions about this subgroup \citep{ebctcg2011}.

\subsection{National Cancer Database Analyses}

The largest study directly addressing endocrine therapy in ER−/PR+ breast cancer was presented at the 2021 ASCO Annual Meeting \citep{asco2021ncdb}. This retrospective analysis of the NCDB examined the impact of adjuvant endocrine therapy on overall survival in ER−/PR+ locoregional breast cancer.

In a propensity-matched cohort of 3,980 patients with ER−/PR+/HER2− disease, endocrine therapy was not significantly associated with improved overall survival (hazard ratio 1.05; 95% confidence interval 0.86–1.28; p = 0.63). This null result suggests that in the HER2-negative context, the PR-positive status alone does not confer endocrine sensitivity sufficient to produce a measurable survival benefit.

Notably, a different result was observed for ER−/PR+/HER2+ disease, where endocrine therapy was associated with improved survival. This differential effect by HER2 status may reflect distinct biology—HER2-positive tumors with PR expression may represent a subgroup with residual luminal signaling or cross-talk between HER2 and hormone receptor pathways that can be therapeutically exploited \citep{asco2021ncdb}.

\subsection{Small Single-Institution Studies}

Several smaller studies have reported conflicting results. A Taiwanese study \citep{cheng2012} of 128 ER−/PR+ patients found that adjuvant tamoxifen improved overall survival and disease-free survival specifically in the lower-grade subgroup (p = 0.01 and p = 0.03, respectively). However, the small sample size, single-institution design, and potential for selection bias limit the generalizability of these findings.

Other single-institution series have reported no significant benefit from endocrine therapy in ER−/PR+ disease, consistent with the NCDB findings and the EBCTCG meta-analysis \citep{kunc2018}.

\subsection{Chemotherapy as Mainstay}

Given the aggressive clinicopathological features of ER−/PR+ tumors and the uncertain benefit of endocrine therapy, chemotherapy has emerged as the preferred systemic treatment modality. Pathological complete response rates to neoadjuvant anthracycline-plus-taxane regimens in ER−/PR+ tumors have been reported as high as 40%, comparable to triple-negative disease \citep{li2023}. This chemosensitivity, combined with the molecular similarity to ER-negative disease, supports chemotherapy as the cornerstone of adjuvant treatment.

\section{Diagnostic Considerations and Retesting}

\subsection{ASCO/CAP Guidelines on ER−/PR+ Results}

The American Society of Clinical Oncology/College of American Pathologists (ASCO/CAP) guidelines on estrogen and progesterone receptor testing explicitly address the ER−/PR+ phenotype \citep{allison2020}. The 2020 guideline update notes that:

PgR may serve as an informal quality control for ER testing. Samples that test ER-negative but PR-positive, especially in the absence of appropriate internal positive controls, should raise suspicion for a false-negative ER result.

Repeat testing in cases with initial ER−/PgR+ results is recommended. The Expert Panel acknowledges that this phenotype is "frequently the result of technical artifact" and supports retesting on the same block, a different tumor block, or using an alternative ER antibody.

Confirmed ER−/PgR+ samples may represent a rare biological phenotype for which endocrine therapy may be offered, though the Panel acknowledges "limited data to support this."

\subsection{Practical Approach to Retesting}

When an ER−/PR+ result is encountered, the following systematic approach is recommended:

First, review pre-analytical variables: cold ischemia time, fixation type and duration, tissue processing conditions. Prolonged cold ischemia (>1 hour) or inadequate fixation can selectively affect ER immunostaining.

Second, assess internal controls: presence of ER-positive normal breast epithelium or other ER-positive elements in the same section can serve as internal controls. Their absence in an ER-negative result should prompt caution.

Third, retest with an alternative antibody clone or on a different tissue block. The two most commonly used ER antibody clones (SP1 and 1D5) may show differential sensitivity in certain contexts.

Fourth, consider the clinical and pathological context: a well-differentiated, low-grade tumor with an ER−/PR+ result should raise stronger suspicion for a false-negative ER than a high-grade tumor.

Fifth, consider molecular profiling: genomic assays such as PAM50, Oncotype DX, or MammaPrint can provide an ER-independent assessment of intrinsic subtype and may help resolve ambiguous cases \citep{prat2012}.

\section{Molecular Profiling and Subtype Reclassification}

\subsection{Gene Expression Profiling}

The advent of gene expression-based classification has provided important insights into the ER−/PR+ phenotype. Using the PAM50 classifier, Prat et al. \citep{prat2012} demonstrated that approximately 25% of clinically ER-negative tumors are classified as luminal by gene expression, suggesting that a subset of ER−/PR+ tumors may indeed harbor luminal biology despite negative ER immunohistochemistry.

This finding has practical implications: patients whose tumors are ER-negative by IHC but luminal by PAM50 may potentially benefit from endocrine therapy, as their gene expression profile suggests underlying hormone receptor pathway activity. Conversely, ER−/PR+ tumors that classify as basal-like by gene expression are unlikely to benefit from endocrine therapy regardless of their PR status.

\subsection{Androgen Receptor Considerations}

The androgen receptor (AR) is expressed in approximately 10–35% of triple-negative breast cancers and defines the luminal androgen receptor (LAR) molecular subtype \citep{lehmann2011}. Some ER−/PR+ tumors may overlap with the LAR subtype, where AR signaling drives tumor growth. While AR-targeted therapies such as enzalutamide and bicalutamide have been explored in AR-positive TNBC with modest results, this represents an alternative endocrine pathway distinct from ER/PR-directed therapy.

\section{Discussion}

\subsection{Synthesizing the Evidence}

The evidence regarding adjuvant endocrine therapy in ER−/PR+/HER2− breast cancer can be summarized as follows:

The EBCTCG meta-analysis, the most robust evidence base available, demonstrates that ER status—not PR status—is the primary predictor of tamoxifen benefit. PR status does not independently predict endocrine therapy efficacy \citep{ebctcg2011}.

The largest database study directly addressing this question (NCDB analysis, n = 3,980 propensity-matched) found no significant overall survival benefit from adjuvant endocrine therapy in ER−/PR+/HER2− disease \citep{asco2021ncdb}.

The biological basis for endocrine therapy benefit in the absence of ER is weak, as the canonical mechanism of endocrine therapy requires ER for therapeutic effect. PR expression in ER-negative tumors may represent technical artifact, ER-independent PR regulation, or cryptic ER activity undetected by IHC.

A subset of ER−/PR+ tumors may harbor luminal gene expression profiles (approximately 25% of clinically ER-negative cases), suggesting potential endocrine sensitivity in a minority of patients \citep{prat2012}.

\subsection{Clinical Implications}

Given this evidence, we propose the following clinical framework for managing ER−/PR+/HER2− breast cancer:

Retesting should be mandatory. All ER−/PR+ results should be retested per ASCO/CAP recommendations. If ER converts to positive on retesting, standard ER-positive treatment guidelines should be followed.

For confirmed ER−/PR+ disease, chemotherapy should be the primary systemic treatment, consistent with treatment approaches for triple-negative breast cancer. Neoadjuvant chemotherapy may be preferred to assess chemosensitivity.

Endocrine therapy cannot be routinely recommended based on current evidence. The NCDB data showing no survival benefit in HER2-negative disease, combined with the EBCTCG finding that PR status does not predict tamoxifen efficacy, do not support routine adjuvant endocrine therapy.

However, endocrine therapy may be considered in selected cases. Patients with confirmed ER−/PR+ tumors that are low-grade, demonstrate luminal gene expression profiles, or have other features suggesting hormone sensitivity (e.g., strong diffuse PR expression) may reasonably be offered endocrine therapy after informed discussion of the limited evidence. Multidisciplinary tumor board discussion is strongly recommended.

Molecular profiling should be encouraged. Gene expression assays can help resolve the biological nature of the tumor and may identify the subset of patients most likely to benefit from endocrine therapy.

Clinical trial enrollment should be prioritized. Given the rarity of this subtype and the paucity of prospective data, patients with confirmed ER−/PR+ disease should be encouraged to enroll in clinical trials when available.

\section{Limitations}

This review has several limitations. The ER−/PR+ phenotype is rare, resulting in small sample sizes in most available studies and limiting statistical power. The majority of evidence comes from retrospective database analyses, which are subject to selection bias, unmeasured confounding, and inaccuracies in diagnostic coding. Prospective randomized trials specifically addressing adjuvant endocrine therapy in ER−/PR+ breast cancer do not exist and are unlikely to be conducted given the rarity of the subtype. Additionally, heterogeneity in IHC methods, antibody clones, scoring criteria, and PR expression thresholds across studies complicates cross-study comparisons.

\section{Conclusion}

Estrogen receptor-negative/progesterone receptor-positive/HER2-negative breast cancer remains one of the most controversial and poorly understood subtypes in breast oncology. The weight of available evidence—from the EBCTCG meta-analysis demonstrating ER as the sole independent predictor of endocrine therapy benefit, to the NCDB propensity-matched analysis showing no survival advantage from endocrine therapy in ER−/PR+/HER2− disease—does not support the routine use of adjuvant endocrine therapy in this population.

However, the heterogeneity of this subtype, the possibility that some tumors harbor genuine luminal biology masked by false-negative ER immunohistochemistry, and the absence of prospective trial data all argue against categorical exclusion of endocrine therapy. A nuanced, individualized approach—incorporating mandatory retesting, molecular profiling, multidisciplinary discussion, and shared decision-making with patients—represents the most evidence-based strategy for managing this rare and challenging breast cancer subtype.

Future research should focus on large-scale multi-institutional prospective studies with standardized diagnostic criteria, mandatory retesting protocols, and integrated genomic profiling to definitively characterize the biology and treatment responsiveness of confirmed ER−/PR+ breast cancer.

\section{Broader Impact}

The clinical challenge posed by ER−/PR+/HER2− breast cancer extends beyond this specific subtype to broader questions in precision oncology: How do we make treatment decisions when the evidence base is sparse? How do we reconcile discordant biomarker results with established treatment paradigms? And how do we integrate increasingly sophisticated molecular profiling tools into a clinical framework traditionally built on immunohistochemistry? Addressing these questions for the ER−/PR+ subtype may provide insights applicable to other rare and ambiguous breast cancer phenotypes, ultimately advancing the goal of truly individualized cancer care.


\bibliographystyle{plainnat}
\bibliography{references}

\end{document}